\documentclass[12pt]{article}
\usepackage[utf8]{inputenc}
\usepackage[T1]{fontenc}
\usepackage{hyperref}
\title{Sengele Language Resources for K-12 Curriculum Development}
\author{Compiled by Claude Sonnet 5.5 (claude-sonnet-5-5) for Andr\'as Kornai}
\date{}
\begin{document}
\maketitle

\section{Introduction}
Sengele (ISO 639-3 szg; also called Kesengele or Sengere) is a Bantu language of the Democratic Republic of the Congo. Wikipedia places it in Bantu Zone C, Bangi--Ntomba (Guthrie C.33), and gives 17,000 speakers (2002). Omniglot gives a similar figure and locates the speakers between the Congo River and Lake Mai-Ndombe. Joshua Project lists 29,000 Sengele people. Dialects named in the sources are Ng\"ong\"o, Bok\'ot\'e, Mp\'eng\'e and Mb\'elo. The language is tonal. Search results turned up very little material for schools, so this report is short.

\section{Primary Educational Resources}
No primers, textbooks or school materials in Sengele turned up in the search results.

\section{Linguistic and Reference Materials}
\begin{itemize}
\item Wikipedia article on the Sengele language: \url{https://en.wikipedia.org/wiki/Sengele_language}
\item Omniglot page on Sengele, covering the writing system and tone marking: \url{https://omniglot.com/writing/sengele.htm}
\item A 1997 study of the Mb\'elo dialect by Motingea Mangulu, which draws on earlier notes by Hulstaert, is mentioned in the Wikipedia summary. A search result on AJOL may relate to it, but I did not confirm that: \url{https://www.ajol.info/index.php/aq/article/view/5596}
\end{itemize}

\section{Religious and Community Resources}
\begin{itemize}
\item Joshua Project profile of the Sengele people, which reports that most community members speak the language at home: \url{https://joshuaproject.net/people_groups/14852}
\item A Sengele version of the JESUS Film is reported in the search summary. I saw no direct URL for it, so it is not in the JSON list.
\end{itemize}

\section{Digital and Audio Resources}
I found no corpora, dictionaries or audio collections online. The JESUS Film version is the only audio-visual item mentioned.

\section{Curriculum Development Context}
The language is said to be used as a first language by the whole community but is not known to be taught in schools. A curriculum would have to start from basics: agreeing an orthography (Omniglot is a starting point), recording speakers, and collecting texts. Local organizations and linguists should be contacted directly. I found none by name.
\end{document}